% FILE: 06_contribution_to_thesis.tex
% Project: livestream
% Thesis Role: real-time-process-broadcast-and-livestream-standards

\section{Contribution to the Dissertation}
\label{sec:livestream-contribution}

\subsection{Contribution to Prediction vs.~Coordination}
\label{sec:livestream-contribution-prediction}

The dissertation's central argument is that prediction and coordination are
categorically distinct: a system that predicts the form of process outcomes does
not thereby coordinate the actual process. The livestream project provides the
most direct operational demonstration of this distinction in the entire foundry.

The five-agent AGI livestreaming swarm does not predict whether the factory is
conforming to its declared process model. It observes, in real time, whether the
observed event sequence is consistent with the declared Petri net model, and it
acts on that observation through the MAPE-K autonomic controller. This is
coordination by process evidence, not prediction by pattern completion. The EWMA
drift detector (alpha\,=\,0.15, LCL\,=\,0.920) ensures that the factory cannot
coast on a historically good conformance score: each new event is weighted into
the moving average, and a sustained conformance drop triggers a remediation
response before the next downstream actuation is authorized.

The conformance audit result (Fitness\,=\,0.9880, Precision\,=\,1.0000) provides
empirical grounding for this claim: the AGI conversation trace, treated as a process
execution, is highly consistent with the Blue River Dam thesis Petri net model.
This means the factory's coordination activity---the sequence of commits, decisions,
and doctrine emissions that constitute the research program---is not arbitrary. It
follows a structured process that can be described, checked, and enforced.

\subsection{Contribution to Process Evidence}
\label{sec:livestream-contribution-evidence}

The livestream project extends the concept of process evidence from static event
logs (collected, analyzed, reported) to live event streams (emitted, checked,
acted upon). The \texttt{aalst\_broadcaster.js} module demonstrates that a Git
repository's commit history can be transformed into an OCEL 2.0-conformant event
stream with minimal instrumentation: a single \texttt{git log} invocation,
a pipe-delimiter parse, and a UTC ISO 8601 timestamp conversion. This establishes
a low-friction path for any software factory to begin emitting process evidence
without requiring changes to its existing toolchain.

The \texttt{master\_ocel\_stream.js} polling loop demonstrates that a filesystem
directory of replay files can serve as a live event source, bridging the gap
between asynchronous file-based data exchange (the zoeapp replay convention) and
synchronous stream-based process monitoring. The one-second polling interval is
a pragmatic choice that trades latency for simplicity; the Stream Director audit
confirms that the rendering latency (6--15\,ms) is dominated by the browser
simulation loop, not the filesystem poll.

The Petri net snapshot (final sound marking, token at \texttt{[Decommissioning]},
no deadlocks) is evidence that the process modeled by the broadcast architecture
is \emph{sound}: it can always reach the terminal state, and it does so with
exactly one token. Soundness is a necessary condition for a workflow net to be
suitable for conformance checking; a non-sound net may produce spurious
conformance results.

\subsection{Contribution to Manufacturing Architecture}
\label{sec:livestream-contribution-manufacturing}

\textbf{[AUTHOR THESIS]} The livestream project contributes to the manufacturing
architecture by demonstrating that the factory's own operation can be the subject
of OCEL 2.0 event logging. In conventional CodeManufactory usage, OCEL 2.0 logs
are emitted by the artifacts being manufactured (e.g., the wasm4pm execution engine
emitting events as it processes an input event log). The livestream project inverts
this: the factory itself emits OCEL 2.0 events as it manufactures artifacts. This
inversion makes the manufacturing pipeline itself auditable by the same process
mining tools it manufactures support for.

The SHA-256 cryptographic ledger chain links every ledger block to its predecessor,
creating a tamper-evident record of the factory's manufacturing history. This
satisfies the foundry's receipt doctrine at the event level: not only are major
research findings certified by named receipts (as in the RECEIPT\_REGISTRY), but
every individual factory event is individually chained and verifiable. The M\&A
Claims Verification Board in the browser visualizer demonstrates the downstream
consequence: a board-admissible M\&A claim backed by a tampered ledger block is
automatically rejected, enforcing the principle that capital-flow claims must be
grounded in verifiable process evidence.

The hot-swappable S-component architecture (\texttt{hotSwapSComponent()}) allows
the declared process model to be updated as the factory evolves, without stopping
the conformance checking loop. This is a necessary feature for a living research
program whose process model evolves as new doctrine is issued.

\subsection{Contribution to Capital-Flow Infrastructure}
\label{sec:livestream-contribution-capital}

The M\&A Claims Verification Board implemented in the browser visualizer represents
the dissertation's most concrete instantiation of the capital-flow infrastructure
argument. Three specific M\&A claim categories are evaluated in real time against
the live event stream:
\begin{itemize}
  \item EBITDA synergy claims (\$5M target): validated against process fitness
    metrics from the conformance stream;
  \item Compliance enforcement claims (\$3.5M target): validated against DECLARE
    constraint satisfaction rates in the LTL monitor;
  \item Data room lineage claims (\$1.2M target): validated against the
    cryptographic ledger chain integrity.
\end{itemize}

This architecture demonstrates that capital-flow claims---which in conventional M\&A
due diligence are supported by static documents and management representations---can
instead be grounded in live process evidence subject to real-time conformance
checking. When the AuditLedger detects a tampered block and automatically rejects
the associated lineage claim, it is demonstrating that process evidence can serve
as an automated gating mechanism for capital allocation decisions.

\textbf{[FUTURE WORK]} The current implementation evaluates synthetic M\&A claims
against simulated traces. Connecting the Verification Board to live enterprise event
streams from actual business processes would constitute a production-grade deployment
of this architecture.
